\scanpage{41}

Alors

a) $M$ est un sous-fibré vectoriel de rang $2$ de $\underline{\operatorname{End}}(E)$, qui est un sous-$A$-Module à gauche et à droite. Les deux structures gauche et droite sont reliées par

\[
u\Theta=\Theta\bar u,\qquad\text{pour }\Theta\in\Gamma M,\ u\in\Gamma A
\]

(elles se déterminent \inserted{donc} mutuellement \ldots) Le $A$-Module $M$ à gauche \deleted{est loc libre} inv. (donc aussi celui à droite).

b) $\Pi$ est formé des sections $\Theta$ de $M$ telles que $N(\Theta)=-1$. Il forme un torseur$^*$ sous $SW(A)\subset W(A)^*$, le \deleted{groupe} schéma en groupes des $x\in\Gamma A$ tels que $N(x)=1$, par l'opération de multiplication à g. (\unsure{aussi bien|ou à droite} \inserted{par celle à} droite).

% slanted
\margin{\uline{NB} Ce torseur est en fait loc. trivial (Zar)}

\deleted{\uline{Dém.} OPS $E=A=\underline O[u]/(u^2+bu+c)$. Alors $W$ est défini comme le noyau de}

\[
\deleted{\unsure{\mathscr L_u}:\underline{\operatorname{End}}(\unsure{.}}
\]

c) Considérons

\[
\underset{\substack{\cup\\X}}{P}=P(E),\qquad P^*=P-\underset{\substack{\Vert\\\operatorname{Spec}(A)}}{X'}
\]

le fibré des droites tot. non isotropes de $A$.

$W(A)^*$ opère sur $P^*$ par son opération sur $E$, par $u\in\Gamma W(A)^*$ définit la transformation $D\longmapsto u.D$ (NB. $\unsure{G_m^{*}|G_m}=W(\underline O)^*\subset A^*$ opère trivialement donc c'est le quotient $W(A)^*/G_m$ qui opère \ldots).

On a un \deleted{\unsure{applicat|application}} morphisme canonique

\diagram{\begin{tikzcd}[ampersand replacement=\&,row sep=1pt]P^*\arrow[r,"f"]\&\Pi\\ D\arrow[r,maps to]\&\sigma_D\end{tikzcd}}

% archive sheet number: 40
